% Source PDF pages 95--99; manual pages 2-66--2-70.
\subsection{Ranges and Distances}\label{sec:ranges-distances}

Projecting a flight path onto the earth's surface produces the ground range described
in \cref{sec:ground-speed}. Ground range measures distance traveled but is
not referenced to a fixed origin or direction. The quantities in this section use
surface coordinate systems to describe a vehicle's ground track relative to a known
reference point and azimuth.

\taossubhead{East and North Distances}

An east/north coordinate system is defined on the earth's surface. Its east axis
follows a line of constant latitude and its north axis follows a line of constant
longitude, as shown in \cref{fig:east-north-distances}. The origin is a known
reference point, commonly the initial vehicle position. This two-dimensional system
is most useful for short trajectories for which earth curvature is not important.

\begin{figure}[htbp]
  \centering
  \resizebox{0.82\linewidth}{!}{\begin{tikzpicture}[>=Latex,line cap=round,line join=round,scale=0.78]
  \foreach \y in {0,1.2,2.4}{\draw[densely dashed] (-3.6,\y)--(4.0,\y);}
  \foreach \x in {-2.4,-1.2,0,1.2,2.4,3.6}{\draw[densely dashed] (\x,-0.5)--(\x,3.0);}
  \coordinate (o) at (-1.2,0.6);
  \coordinate (v) at (2.4,2.4);
  \fill (o) circle (2.4pt) node[below left,align=center] {reference\\point};
  \fill (v) circle (2.4pt) node[right] {vehicle};
  \draw[->,very thick] (o)--(3.2,0.6) node[midway,below] {east axis};
  \draw[->,very thick] (o)--(-1.2,2.9) node[above] {north axis};
  \draw[<->,thick] (-1.2,2.4)--(2.4,2.4) node[midway,above] {east};
  \draw[<->,thick] (2.4,0.6)--(2.4,2.4) node[midway,right] {north};
  \node[align=left] at (-4.6,2.4) {lines of\\constant latitude};
  \node at (0.6,-0.85) {lines of constant longitude};
\end{tikzpicture}
}
  \caption{East and north distances.}
  \label{fig:east-north-distances}
\end{figure}

The east distance is measured between the reference longitude and the vehicle
longitude along the reference latitude. The north distance is measured between the
reference latitude and vehicle latitude along the reference longitude. Both
calculations require the surface distance between two known longitude/latitude
points.

\taossubhead{Sodano's Inverse Method}

TAOS uses Sodano's inverse method because it is noniterative and accurate over both
short and long distances.\manualref{16} The function \texttt{sodano\_inverse} in
the \texttt{derivs} module takes two surface points
$(\lambda_1,\delta_{gd_1})$ and $(\lambda_2,\delta_{gd_2})$. Their ellipsoidal
surface distance $\Delta r_s$ is
\begin{equation}
\begin{aligned}
\frac{\Delta r_s}{R_p}={}&(1+f+f^2)\phi
 +\sin\beta_1\sin\beta_2
  \left[(f+f^2)\sin\phi-\frac12 f^2\phi^2\csc\phi\right]\\
&-\frac12m\left[(f+f^2)(\phi+\sin\phi\cos\phi)
                 -f^2\phi^2\cot\phi\right]\\
&-\frac12\sin^2\beta_1\sin^2\beta_2\,f^2\sin\phi\cos\phi\\
&+\frac12\sin\beta_1\sin\beta_2\,m f^2
  \left[\sin\phi\cos^2\phi+\phi^2\csc\phi\right]\\
&+\frac1{16}f^2m^2
  \left[\phi+\sin\phi\cos\phi-2\sin\phi\cos^3\phi
        -8\phi^2\cot\phi\right].
\end{aligned}
\label{eq:sodano-inverse-distance}
\end{equation}
Here the auxiliary central angle $\phi$ and quantity $m$ are
\begin{equation}
  \cos\phi
  =\sin\beta_1\sin\beta_2
   +\cos\beta_1\cos\beta_2\cos(\lambda_2-\lambda_1),
  \label{eq:sodano-central-angle}
\end{equation}
\begin{equation}
  m=1-\cos^2\beta_1\cos^2\beta_2
       \sin^2(\lambda_2-\lambda_1)\csc^2\phi,
  \label{eq:sodano-inverse-m}
\end{equation}
and the reduced latitudes are
\begin{equation}
  \tan\beta_1
  =(1-f)\tan\delta_{gd_1}
  =\frac{R_p\sin\delta_{gd_1}}{R_{\oplus}\cos\delta_{gd_1}},
  \label{eq:sodano-beta1}
\end{equation}
\begin{equation}
  \tan\beta_2
  =(1-f)\tan\delta_{gd_2}
  =\frac{R_p\sin\delta_{gd_2}}{R_{\oplus}\cos\delta_{gd_2}}.
  \label{eq:sodano-beta2}
\end{equation}
The equatorial radius $R_{\oplus}$, polar radius $R_p$, and flattening $f$ are
defined in \cref{sec:local-geodetic}.

\noindent\textit{Editorial note.} The printed source omits the cosine function
before $(\lambda_2-\lambda_1)$ in \cref{eq:sodano-central-angle}; it is restored
here because the surrounding equations and the spherical law of cosines require it.

Sodano's inverse method also provides the forward and backward azimuths of the
surface curve joining the points. Azimuth is measured clockwise from north, so
$0^\circ$ points north and $90^\circ$ points east. The forward azimuth is evaluated
at the first point and the backward azimuth at the second point:
\begin{equation}
  \tan\psi_{\mathrm{fwd}}
  =\frac{\cos\beta_2\sin\Lambda}
         {\sin\Delta\beta
          +2\cos\beta_2\sin\beta_1\sin^2(\Lambda/2)},
  \label{eq:sodano-forward-azimuth}
\end{equation}
\begin{equation}
  \tan\psi_{\mathrm{back}}
  =\frac{\cos\beta_1\sin\Lambda}
         {\sin\Delta\beta
          -2\cos\beta_1\sin\beta_2\sin^2(\Lambda/2)}.
  \label{eq:sodano-backward-azimuth}
\end{equation}
The corrected longitude difference $\Lambda$ is
\begin{equation}
\begin{aligned}
\Lambda={}&\lambda_2-\lambda_1
 +\frac{\cos\beta_1\cos\beta_2\sin(\lambda_2-\lambda_1)}{\sin\phi}
 \Bigg[(f+f^2)\phi\\
&\quad-\frac12\sin\beta_1\sin\beta_2\,f^2
       \left(\sin\phi+2\phi^2\csc\phi\right)\\
&\quad+\frac14m f^2
       \left(\sin\phi\cos\phi-5\phi+4\phi^2\cot\phi\right)
 \Bigg],
\end{aligned}
\label{eq:sodano-corrected-longitude-difference}
\end{equation}
while the reduced-latitude difference can be written as the series
\begin{equation}
\begin{aligned}
\Delta\beta={}&\delta_{gd_2}-\delta_{gd_1}
+n\left(\sin2\delta_{gd_1}-\sin2\delta_{gd_2}\right)\\
&-\frac12n^2\left(\sin4\delta_{gd_1}-\sin4\delta_{gd_2}\right)
+\frac13n^3\left(\sin6\delta_{gd_1}-\sin6\delta_{gd_2}\right),
\end{aligned}
\label{eq:sodano-delta-beta-series}
\end{equation}
or equivalently
\begin{equation}
\begin{aligned}
\Delta\beta={}&\delta_{gd_2}-\delta_{gd_1}
+2\sin(\delta_{gd_2}-\delta_{gd_1})
 \Big[\sin\beta_1\sin\beta_2(n+n^2+n^3)\\
&\hspace{11em}-\cos\beta_1\cos\beta_2(n-n^2+n^3)\Big],
\end{aligned}
\label{eq:sodano-delta-beta-alternative}
\end{equation}
where
\begin{equation}
  n=\frac{R_{\oplus}-R_p}{R_{\oplus}+R_p}.
  \label{eq:sodano-flattening-n}
\end{equation}

\taossubhead{Downrange and Crossrange}

A second surface coordinate system measures downrange from a reference point along
a reference azimuth and crossrange perpendicular to that azimuth, as shown in
\cref{fig:downrange-crossrange}.

\begin{figure}[htbp]
  \centering
  \resizebox{0.78\linewidth}{!}{\begin{tikzpicture}[>=Latex,line cap=round,line join=round,scale=0.9]
  \coordinate (o) at (-3.5,-1.6);
  \coordinate (p) at (1.0,0.2);
  \coordinate (v) at (0.4,2.0);
  \draw[densely dashed] (o)--(4.2,-1.6) node[right] {latitude};
  \draw[densely dashed] (o)--(-3.5,2.5) node[above] {longitude};
  \draw[->,very thick] (o)--(3.4,1.25) node[pos=0.56,below] {downrange}
    node[right,align=left] {reference\\azimuth};
  \fill (o) circle (2.3pt) node[below left,align=center] {reference\\point};
  \fill (p) circle (2.3pt) node[below right] {P};
  \fill (v) circle (2.3pt) node[above] {vehicle};
  \draw[->,thick] (p)--(v) node[midway,right] {crossrange};
  \draw[->] ($(p)+(22:0.52)$) arc[start angle=22,end angle=108,radius=0.52]
    node[midway,above right] {$\theta$};
  \draw[->] (-3.0,-0.7) arc[start angle=80,end angle=25,radius=0.75]
    node[midway,right,align=center] {azimuth\\angle};
\end{tikzpicture}
}
  \caption{Downrange and crossrange.}
  \label{fig:downrange-crossrange}
\end{figure}

A Newton--Raphson iteration varies the downrange location of point P until the
backward azimuth from P to the reference point and the forward azimuth from P to
the vehicle differ by $90^\circ$:
\begin{equation}
  \theta=\psi_{\mathrm{back}}-\psi_{\mathrm{fwd}}=90^\circ.
  \label{eq:downrange-crossrange-perpendicular}
\end{equation}
The generic Newton--Raphson update is
\begin{equation}
  x_{n+1}
  =x_n-f(x_n)\left[\frac{\partial f(x_n)}{\partial x}\right]^{-1},
  \label{eq:downrange-newton-raphson}
\end{equation}
where $n$ is the iteration count. TAOS estimates the derivative with a forward
difference. For the downrange search, $x$ is the downrange distance and the error
function is
\begin{equation}
  f(x)=\frac{\pi}{2}
       -\left|\psi_{\mathrm{back}}-\psi_{\mathrm{fwd}}\right|.
  \label{eq:downrange-error-function}
\end{equation}

\taossubhead{Sodano's Direct Method}

The error evaluation requires Sodano's direct method. Given reference longitude
$\lambda_1$, reference latitude $\delta_{gd_1}$, reference azimuth $\psi$, and
surface distance $\Delta r_s$, the method computes the backward azimuth,
longitude $\lambda_2$, and latitude $\delta_{gd_2}$ of point P.\manualref{16}
It begins by evaluating
\begin{equation}
\begin{aligned}
\phi={}&\xi-ne'^2\sin\xi
 +\frac12me'^2(\sin\xi\cos\xi-\xi)
 +\frac52n^2e'^4\sin\xi\cos\xi\\
&+\frac1{16}m^2e'^4
 \left(11\xi-13\sin\xi\cos\xi-8\xi\cos^2\xi
       +10\sin\xi\cos^3\xi\right)\\
&+\frac12mne'^4
 \left(3\sin\xi+2\xi\cos\xi-5\sin\xi\cos^2\xi\right),
\end{aligned}
\label{eq:sodano-direct-phi}
\end{equation}
where, for this direct-method derivation,
\begin{equation}
  m=\frac12
    \left(1+\frac12e'^2\sin^2\beta_1\right)
    \left(1-\cos^2\beta_1\sin^2\psi\right),
  \label{eq:sodano-direct-m}
\end{equation}
\begin{equation}
  n=\frac12
    \left(1+\frac12e'^2\sin^2\beta_1\right)
    \left(\sin^2\beta_1\cos\xi
          +\cos\beta_1\cos\psi\sin\beta_1\sin\xi\right),
  \label{eq:sodano-direct-n}
\end{equation}
\begin{equation}
  \xi=\frac{\Delta r_s}{R_p},
  \label{eq:sodano-direct-xi}
\end{equation}
and the source defines
\begin{equation}
  e'=\frac{1}{\sqrt{1-e^2}}
     =\sqrt{\frac{R_{\oplus}^2+R_p^2}{R_p^2}}.
  \label{eq:sodano-eprime}
\end{equation}
The reduced latitude $\beta_1$ is computed from
\cref{eq:sodano-beta1}. The geodetic latitude of P then follows from
\begin{equation}
  \tan\delta_{gd_2}
  =\frac{\tan\beta_2}{1-f}
  =\frac{R_{\oplus}\sin\beta_2}{R_p\cos\beta_2},
  \label{eq:sodano-direct-latitude}
\end{equation}
with
\begin{equation}
  \sin\beta_2
  =\sin\beta_1\cos\phi+\cos\beta_1\cos\psi\sin\phi,
  \label{eq:sodano-direct-sin-beta2}
\end{equation}
\begin{equation}
  \cos\beta_2
  =\sqrt{\cos^2\beta_1\sin^2\psi
          +\left(\sin\beta_1\sin\phi
                 -\cos\beta_1\cos\psi\cos\phi\right)^2}.
  \label{eq:sodano-direct-cos-beta2}
\end{equation}
The longitude is
\begin{equation}
  \lambda_2=\lambda_1+l
  +\tan^{-1}\!\left[
    \frac{\sin\phi\sin\psi}
         {\cos\beta_1\cos\phi-\sin\beta_1\sin\phi\cos\psi}
  \right],
  \label{eq:sodano-direct-longitude}
\end{equation}
where
\begin{equation}
  l=\cos\beta_1\sin\psi
    \left[-f\xi+3f^2n\sin\xi
          +\frac32f^2m(\xi-\sin\xi\cos\xi)\right].
  \label{eq:sodano-direct-l-correction}
\end{equation}
Finally, the backward azimuth from P to the reference point is
\begin{equation}
  \tan\psi_{\mathrm{back}}
  =\frac{-\cos\beta_1\sin\psi}
         {\sin\beta_1\sin\phi-\cos\beta_1\cos\psi\cos\phi}.
  \label{eq:sodano-direct-backward-azimuth}
\end{equation}
The function \texttt{sodano\_direct} implements these equations. Once P is known,
Sodano's inverse method gives the forward azimuth and distance from P to the
vehicle, and \cref{eq:downrange-error-function} supplies the Newton--Raphson
error.

\noindent\textit{Editorial note.} Equation~\ref{eq:sodano-eprime} is reproduced
as printed. Its two equalities are not algebraically consistent with the eccentricity
definition in \cref{eq:ellipsoid-parameter-relation}; the source issue is recorded in
the Chapter 2 editorial metadata rather than silently corrected.
